\documentclass[12pt]{amsart}
\usepackage{amsmath,amsthm,amssymb}
\usepackage{dsfont,bm,mathtools}
\usepackage[usenames,dvipsnames,svgnames,table]{xcolor}
\usepackage[top=2.3cm,bottom=2.3cm,left=2cm,right=2cm]{geometry}
\usepackage{hyperref}
\hypersetup{hidelinks}
\allowdisplaybreaks
\emergencystretch=1.5em
\definecolor{ProposalOrange}{RGB}{160,65,0}
\definecolor{RefereePurple}{RGB}{103,34,147}
\newtheorem{lem}{Lemma}[section]
\theoremstyle{remark}
\newtheorem{rem}[lem]{Remark}
\theoremstyle{definition}
\newtheorem{nota}[lem]{Notation}
\definecolor{ResolutionGreen}{RGB}{0,100,55}
\newcommand{\C}{\mathds{C}}
\newcommand{\Q}{\mathds{Q}}
\newcommand{\rv}{\operatorname{rv}}
\newcommand{\Gal}{\operatorname{Gal}}
\newcommand{\Aut}{\operatorname{Aut}}
\newcommand{\RV}{\mathrm{RV}}
\newcommand{\RES}{\mathrm{RES}}
\newcommand{\VF}{\mathrm{VF}}
\newcommand{\K}{\mathds{k}}
\newcommand{\Z}{\mathds{Z}}
\newcommand{\N}{\mathds{N}}
\newcommand{\vrv}{\operatorname{vrv}}
\newcommand{\GL}{\operatorname{GL}}
\newcommand{\dcl}{\operatorname{dcl}}
\newcommand{\pr}{\operatorname{pr}}
\newcommand{\defn}{\mathrm{def}}
\newcommand{\given}{\,\middle|\,}
\newcommand{\Sbase}{\mathbb S}
\newcommand{\RVlang}{\mathcal L_{\mathrm{RV}}}
\numberwithin{equation}{section}
\pagestyle{plain}

\newcommand{\A}{\mathds{A}}
\newcommand{\Gm}{\mathbb G_m}
\newcommand{\Var}{\mathrm{Var}}
\newcommand{\Kres}{K_0(\RES)}
\newcommand{\Bangres}{{!K_0(\RES)}}
\newcommand{\Kraw}{K_0^{\widehat\mu,\mathrm{raw}}(\Var_{\C})}
\newcommand{\Kmon}{K_0^{\widehat\mu}(\Var_{\C})}
\newcommand{\E}{\mathbb E_b}
\newcommand{\Vol}{\operatorname{Vol}}
\newcommand{\vv}{\operatorname{val}}
\newcommand{\LiftingTheorem}{the lifting theorem of Section~3}

\theoremstyle{plain}
\newtheorem*{thm*}{Theorem}
\begin{document}
\setcounter{section}{5}
\setcounter{subsection}{1}
\setcounter{equation}{1}
{\small\noindent Forey--Yin, \emph{Bounded integral and motivic Milnor fiber}.
Section 5 review: R084--R085.
Approved context is black; proposed manuscript text is orange.}
\medskip

\noindent We work in the original language over
$F=\C((t))$, with value group $\Q$ in the Puiseux field.
The preceding compactness argument gives a definable family
$\mathcal U=\{(\gamma,u):u\in U_\gamma\}$ of graded RV
representatives, with uniform bijections
$H_\gamma:A_\gamma\to\mathbb L\bm U_\gamma$.
All formulas and their witnesses were obtained in a sufficiently
saturated extension before restricting to rational parameters.
Recall
\[
 A^{\mathrm{inf}}
 =\{(x,a):a\in A_{\vv(x)}\},\qquad
 U^{\mathrm{inf}}
 =\{(z,u):u\in U_{\vrv(z)}\},
\]
where $x\in(\VF^\times)^n$ and $z\in(\RV^\times)^n$.
Put $R=\Bangres$ and $g=[\Gm]\in R$.

{\color{ProposalOrange}
\begin{thm*}[Inflated-family Fubini]
There is a finite $F$-definable partition $(E_a)_a$ of $\Q^n$
and classes $v_a\in R$ such that
\begin{equation}\label{trop:fub:cons}
 \nu(\gamma):=\E\!\left(\int[A_\gamma]\right)=v_a
 \quad\text{for }\gamma\in E_a.
\end{equation}
The partition includes the parameters with empty fibers, where
the value is $0$. Moreover,
\begin{equation}\label{fubini:inflated:RES}
 \E\!\left(\int[A^{\mathrm{inf}}]\right)
   =g^n\sum_a\chi_b(E_a)v_a\quad\text{in }R.
\end{equation}
Applying the approved comparison $\Theta$ gives
\[
 \Vol([A^{\mathrm{inf}}])
   =[\Gm]^n\sum_a\chi_b(E_a)\Theta(v_a).
\]
\end{thm*}

\begin{proof}
\emph{A relative monomial decomposition.}
Work first on a fixed zero-coordinate stratum of one tagged
component, and write
$\mathcal U\subseteq\Gamma^n\times(\RV^\times)^r$.
We claim that it has a finite definable partition
$(\mathcal U_i)_i$, integer matrices $N_i\in\GL_r(\Z)$,
fixed nonempty RES sets $W_i$ with singleton value support,
and definable sets
$J_i\subseteq\Gamma^n\times\Gamma^{\ell_i}$, such that
\begin{equation}\label{relative:monomial:chart}
 h_{N_i}(U_{i,\gamma})
       =W_i\times(J_{i,\gamma})^\sharp,\qquad
 J_{i,\gamma}=\{\eta:(\gamma,\eta)\in J_i\}.
\end{equation}
All these data are over $F$, and the coordinate changes act on
$u$ alone, leaving the parameter $\gamma$ unchanged.
The finite labels are transported with the charts as in
Section~3. Combine the finitely many tag and zero strata at
the end.

Here is a proof of the claim from Lemma~5.12.
Apply that lemma to $X=U^{\mathrm{inf}}$, in coordinates
$(z,u)$, to get value-induced pieces
\[
 X_i=X\cap\vrv^{-1}(D_i),\qquad
 h_M(X_i)=V\times I^\sharp,\qquad \vrv(V)=\{\beta\}.
\]
Fix one such piece and let $m$ be the number of coordinates of
$V$. Each $X_i$ is invariant under multiplication of $z$ by
$(\K^\times)^n$: both the family condition and $D_i$
depend on $z$ only through its values.

Let $B$ be the block of $M$ in its first $m$ rows and first
$n$ columns, and put $p=\operatorname{rank}_{\Q}B$,
$h=m-p$. The induced action on $V$ has exponent matrix $B$.
Smith reduction gives $Q\in\GL_m(\Z)$ such that, after this
change of active coordinates, its image is
\[
 \{1\}^h\times(\K^\times)^p.
\]
Indeed the nonzero diagonal powers in Smith form are
surjective on the units of the algebraically closed residue
field. Since the remaining factor $I^\sharp$ already contains
every angular coordinate, invariance of $X_i$ implies
invariance of $V$ under this image. It follows that
\[
 h_Q(V)=W_i\times(\beta_{\mathrm{free}})^\sharp
\]
for a definable $W_i$ in the first $h$ fixed-value
coordinates.

The first $h$ rows of $\operatorname{diag}(Q,1)M$ vanish in
the $z$ columns. Their restrictions to the $u$ columns form
a primitive sublattice of $\Z^r$: their span is primitive
already in $\Z^{n+r}$, since these rows belong to an integer
basis. Extend these rows to a basis of $\Z^r$, obtaining
$N_i\in\GL_r(\Z)$. Write
\[
 h_{N_i}(u)=(w,y),\qquad
 \vrv(W_i)=\{\beta_i^0\},\qquad \ell_i=r-h.
\]
On a fixed joint value tuple, membership in $X_i$ is now
equivalent to $w\in W_i$. All other angular coordinates
are free by the preceding invariance argument.
Consequently the required value set is
\[
 J_i=\{(\gamma,\eta):
       (\gamma,N_i^{-1}(\beta_i^0,\eta))\in D_i\}.
\]
This proves (\ref{relative:monomial:chart}).
It also proves original-base definability: $W_i$ is a
definable projection, its value tuple is fixed over $F$,
and all coordinate matrices are integral. No angular
representative is retained as a parameter.

\emph{A separate base partition.}
The projections of the value pieces to the parameter space
may overlap. For example, in the constant family
$U_\gamma=\{1,2\rv(t),3\rv(t)\}$, the natural pieces at
internal values $0$ and $1$ both project onto the whole base.
We therefore index the source pieces by $i$ and construct
the base pieces independently.

For each $i$, put $e_i(\gamma)=\chi_b(J_{i,\gamma})$.
These are finite-range constructible integer functions.
For completeness, quantifier elimination and
inclusion--exclusion express the indicator of a definable
subset of $\Q^{n+\ell}$ as a finite integer combination of
indicators of closed rational polyhedra.
Every nonempty closed polyhedron has bounded Euler
characteristic $1$: its sufficiently large box cuts are
nonempty closed bounded polyhedra, with geometric Euler
characteristic $1$. A projection and a nonempty fiber of
a closed polyhedron are closed polyhedra again.
Thus for its indicator the fiber integral is the indicator
of the projection. Linearity proves finite-range
constructibility and the projection identity
\begin{equation}\label{gamma:projection:formula}
 \chi_b(J_i)=\int_{\Q^n} e_i(\gamma)\,d\chi_b.
\end{equation}

Take the common finite Boolean refinement of all the supports
$\operatorname{pr}_{\Gamma^n}(J_i)$ and the level sets of
all $e_i$. This is the partition $(E_a)_a$.
It includes the complement of the union of those supports.
Write $e_{ia}$ for the value of $e_i$ on $E_a$, and set
\[
 c_i=[W_i]g^{\ell_i},\qquad v_a=\sum_i e_{ia}c_i\quad\text{in }R.
\]
For every rational $\gamma$, all its coordinates are
$F$-definable. The fiber representatives and
(\ref{relative:monomial:chart}), evaluated by the retraction
of Section~4, give
\[
 \E\!\left(\int[A_\gamma]\right)
       =\sum_i\chi_b(J_{i,\gamma})[W_i]g^{\ell_i}.
\]
This proves (\ref{trop:fub:cons}), including empty fibers.
Every contribution indexed by $i$ is constant on each $E_a$;
their projections themselves are not required to be a partition.

\newpage
\emph{The total inflated space.}
Give $U^{\mathrm{inf}}$ the label
$(z,f_{\vrv(z)}(u))$, where $f_\gamma$ is the retained
label of $\bm U_\gamma$; keep the finite tags and degrees.
Its lift is definably bijective to $A^{\mathrm{inf}}$:
use $H_{\vv(x)}$ on the second factor and retain $x$.
The base RV tag is then $z=\rv(x)$. Thus this replacement
is justified by an actual uniform bijection.

The relative coordinate changes send the corresponding pieces
of $U^{\mathrm{inf}}$ to
\[
 W_i\times J_i^\sharp,
\]
after a coordinate permutation. Here $J_i^\sharp$ is the
full RV lift in $n+\ell_i$ coordinates.
The Section~4 product formula therefore gives directly
\[
 \E\!\left(\int[A^{\mathrm{inf}}]\right)
   =\sum_i\chi_b(J_i)[W_i]g^{n+\ell_i}
   =g^n\sum_i\chi_b(J_i)c_i.
\]
Using (\ref{gamma:projection:formula}) and the base partition,
\[
 \sum_i\chi_b(J_i)c_i
     =\sum_a\chi_b(E_a)\sum_i e_{ia}c_i
     =\sum_a\chi_b(E_a)v_a.
\]
This proves (\ref{fubini:inflated:RES}). Applying $\Theta$
proves the displayed motivic-volume identity.
\end{proof}

The factor $[\Gm]^n$ comes from the additional $n$ VF value
coordinates in $A^{\mathrm{inf}}$.
For $A_\gamma=B\cap\vv^{-1}(\gamma)$, the formula above still
concerns this inflated total space.
An identity for the original $B$ requires a decomposition
compatible with its own tropicalization; it does not follow by
cancelling $[\Gm]^n$ in the Grothendieck ring.
}
\medskip

{\normalfont\small\color{RefereePurple}
\noindent\textbf{Referee R084 [Error] -- Projected pieces need not
form a partition.}
Use an independently indexed base partition, sum all source-piece
contributions over it, and include empty fibers.\par\medskip
\noindent\textbf{Referee R085 [Clarification; revised] --
The extra torus factor.}
The stated total space is an inflated family. Its extra torus
factor cannot be cancelled to obtain the usual tropical Fubini
identity for an uninflated set.\par}
\medskip

{\normalfont\small\color{ProposalOrange}
\noindent\textbf{Proposed disposition of R084--R085.}
Replace the old three bullets and the following computation
by the statement and proof above. The relative decomposition is
proved from the approved Lemma~5.12; it supplies the fixed
RES coefficients needed for the separate base partition.
The total identity is first proved in the original RES target,
then transported by $\Theta$.
The elementary polyhedral projection calculation also appears
in \href{https://web.ma.utexas.edu/users/sampayne/pdf/TropicalFubini.pdf}{Nicaise--Payne,
Proposition 3.1.2}; their Theorem 3.1.3 concerns the
uninflated tropicalization identity. Neither result is used
as a proof input here.}
\end{document}
